\subsection{qsection.txt}
%
\subsubsection{Purpose}
This file stores the instantaneous discharge crossing each qsection every \param{dt\_time\_series}. 
For each Q Section, the stored discharge is the sum of the discharge crossing each of its edges.
When an edge is defined in the input file from its south node to its north node, a positive discharge crosses the edge from west to east.
%
\subsubsection{Format description}
\texttt{<Tab>}: tabulation character. \\
    \begin{tabular}{p{6cm} p{9cm}}
    \hline
    \texttt{time <Tab> name\_1 <Tab> \ldots <Tab> name\_N} & This is the header line. \texttt{name\_1}, \ldots, \texttt{name\_N} represent the q section names as defined in the \file{input/qsection.txt} file \\
    \hline 
    $t$ \texttt{<Tab>} $Q_1$ \texttt{<Tab>} \ldots \texttt{<Tab>} $Q_N$ & $t$ is the stored time, $Q_i$ is the discharge (in $\mathrm{m^3.s^{-1}}$) through the q section \texttt{name\_i} \\
    \hline
    \ldots & one line per stored time \\ \hline
    \end{tabular}
%
\subsubsection{Example}
\begin{verbatim}
time	qsection_1	qsection_2	qsection_3
2.0000000000e+1	0.0000000000e+0	4.3274586317e+1	-1.4193271476e+1
3.0000000000e+1	0.0000000000e+0	4.3149535653e+1	-1.4227591312e+1
4.0000000000e+1	0.0000000000e+0	4.3074013841e+1	-1.4246148173e+1
5.0000000000e+1	0.0000000000e+0	4.3030038602e+1	-1.4258119344e+1
6.0000000000e+1	0.0000000000e+0	4.3000513361e+1	-1.4266604580e+1
\end{verbatim}